\documentclass[a4paper,11pt]{article}
\usepackage[T1]{fontenc}
\usepackage[utf8]{inputenc}
\usepackage[french]{babel}
\usepackage{geometry}
\geometry{top=1.5cm, bottom=1.8cm, left=1.8cm, right=1.8cm}
\usepackage{xcolor}
\renewcommand{\familydefault}{\sfdefault}
\definecolor{Couleur1}{HTML}{8b1f30}
\definecolor{Couleur6}{HTML}{CF6876}
\definecolor{Couleur8}{HTML}{D36740}
\usepackage{amsmath}
\usepackage{array, colortbl}
\usepackage{tikz}
\usetikzlibrary{calc}
\usepackage{fontawesome5}
\usepackage{mdframed}
\usepackage{fancyhdr}
\usepackage{enumitem}
\usepackage{multicol}
\usepackage{rotating}
\setlength{\parindent}{0pt}
\newcommand{\lignepoints}{\noindent\dotfill\vspace{0.4em}\par}
\newcommand{\coeur}{\textcolor{Couleur8}{\faHeart}\ }
\newcommand{\coeurrappel}{\textcolor{Couleur6}{\faHeart}\ }

\pagestyle{fancy}
\fancyhf{}
\renewcommand{\headrulewidth}{0pt}
\fancyfoot[L]{\small Mme Bertail}
\fancyfoot[C]{\small \textcolor{Couleur6}{Interrogation de leçon n°2 — Chapitres 1 \& 2}}
\fancyfoot[R]{\small \thepage}

\begin{document}

\begin{mdframed}[linecolor=Couleur1,linewidth=2pt,topline=true,bottomline=true,
  leftline=false,rightline=false,innertopmargin=8pt,innerbottommargin=8pt,
  backgroundcolor=Couleur1!5]
\begin{center}
  {\Large\bfseries\textcolor{Couleur1}{Interrogation de leçon n°2}}\\[4pt]
  {\large\textcolor{Couleur8}{Chapitres 1 \& 2 — Nombres entiers · Droites et segments}}
\end{center}
\end{mdframed}

\vspace{0.4cm}
\noindent\textbf{Nom :} \dotfill\quad\textbf{Prénom :} \dotfill\quad
\textbf{Classe :} \dotfill\quad\textbf{Date :} \dotfill
\vspace{0.5cm}

%--- Séparateur Ch 1 ---
\noindent{\textcolor{Couleur8}{\rule{\textwidth}{0.5pt}}}\\
{\small\textit{\textcolor{Couleur8}{Rappel — Chapitre 1}}}
\vspace{0.1cm}

%--- Q1 ---
\noindent\coeurrappel \textbf{Placer le nombre $2\,360\,174\,085$ dans le tableau ci-dessous.}

\vspace{0.3cm}
\begin{center}
\arrayrulecolor{Couleur1}
\renewcommand{\arraystretch}{1.5}
\scalebox{0.82}{%
\begin{tabular}{|*{12}{>{\centering\arraybackslash}p{1.15cm}|}}
\hline
\rowcolor{Couleur1!10}
\multicolumn{3}{|c|}{\textcolor{Couleur1}{\textbf{Classe des milliards}}} &
\multicolumn{3}{c|}{\textcolor{Couleur1}{\textbf{Classe des millions}}} &
\multicolumn{3}{c|}{\textcolor{Couleur1}{\textbf{Classe des milliers}}} &
\multicolumn{3}{c|}{\textcolor{Couleur1}{\textbf{Classe des unités}}} \\
\hline
 & & & & & & & & & & & \\[1.1cm]
\hline
\end{tabular}
}
\end{center}

\vspace{0.3cm}

%--- Séparateur Ch 2 ---
\noindent{\textcolor{Couleur8}{\rule{\textwidth}{0.5pt}}}\\
{\small\textit{\textcolor{Couleur8}{Nouveau — Chapitre 2}}}
\vspace{0.1cm}

%--- Q2 ---
\noindent\coeur \textbf{Compléter la définition.}

\vspace{0.2cm}
\begin{mdframed}[linecolor=Couleur6,linewidth=4pt,topline=false,bottomline=false,
  rightline=false,backgroundcolor=Couleur6!10,innerleftmargin=8pt,
  innertopmargin=6pt,innerbottommargin=6pt]
On dit que des points sont \textbf{alignés} lorsqu'ils \dotfill
\end{mdframed}

\vspace{0.4cm}

%--- Q3 ---
\noindent\coeur \textbf{Sur la figure ci-dessous, répondre aux questions.}

\vspace{0.3cm}
\begin{minipage}[t]{0.52\textwidth}
\vspace{0pt}
\begin{enumerate}[label=\textcolor{Couleur8}{\textbf{\arabic*.}}, leftmargin=1.5em, itemsep=0.6em]
  \item Donner deux noms différents de la droite passant par A et C.
  \lignepoints
  \lignepoints\lignepoints
  \item Compléter par $\in$ ou $\notin$ :
  \begin{multicols}{2}
  \begin{itemize}[nosep, label=\textcolor{Couleur6}{$\bullet$}]
    \item $A \ldots (d)$
    \item $B \ldots (d)$
    \item $C \ldots [AB]$
    \item $D \ldots [AC)$
  \end{itemize}
  \end{multicols}
\end{enumerate}
\end{minipage}
\hfill
\begin{minipage}[t]{0.44\textwidth}
\vspace{0pt}
\begin{center}
\begin{tikzpicture}[scale=0.9]
  % Droite (d) oblique
  \draw[thick] (-0.5,-0.85) -- (3.8,1.85);
  \node[right] at (3.6,1.5) {$(d)$};

  % Points sur (d)
  \coordinate (A) at (0.3,-0.4);
  \coordinate (C) at (2.4,1.0);
  % Point hors droite
  \coordinate (B) at (1.5,2.0);
  % Point D hors droite
  \coordinate (D) at (3.2,-0.3);

  \foreach \p/\lbl/\pos in {A/A/below left, C/C/above right, B/B/above, D/D/below right}{
    \draw[thick] ($(\p)+(-0.08,-0.08)$) -- ($(\p)+(0.08,0.08)$);
    \draw[thick] ($(\p)+(-0.08,0.08)$) -- ($(\p)+(0.08,-0.08)$);
    \node[\pos] at (\p) {\lbl};
  }
\end{tikzpicture}
\end{center}
\end{minipage}

\end{document}
